% Anonymous English manuscript skeleton; do not add the final option.
% Replace each \PaperTODO with reviewed content. The paper self-check rejects TODOs.
\documentclass[english]{ai4science_ysda}

\newcommand{\PaperTODO}[1]{\textbf{[Draft: #1]}}
\TITLE{Gold Nanoparticle Size in Pulsed Laser Ablation in Liquids: Claims and Data}
\ShortTITLE{Gold Nanoparticle Size: Claims and Data}
\Abstract{\PaperTODO{State the question, method, main finding, and its limits.}}
\KEYWORDS{pulsed laser ablation in liquids, gold nanoparticles, fluence, pulse energy}

% Generate results.tex from the finalized result files before writing results.
\InputIfFileExists{results.tex}{}{}

\begin{document}

\section*{Introduction}
\generated[revised from the first draft against WRITING.md and research/problem.md]{}

Pulsed laser ablation in liquids (LAL) produces colloidal nanoparticles by
focusing laser pulses on a solid target immersed in a liquid. The target, the
liquid and its solutes, the wavelength, pulse duration, pulse energy,
repetition rate and focusing geometry can each be varied, and the combinations
have been used to prepare nanomaterials from much of the periodic
table~\cite{zhang2017}. The freedom is bounded: monomodal, monodisperse
colloids remain difficult to obtain by varying the laser parameters alone, and
the parameters act together, their mutual cross-effects making a general
concept difficult to develop~\cite{zhang2017}.

The process passes through a plasma plume, a shock wave and a cavitation bubble
within nanoseconds, at temperatures of thousands of kelvins and under steep
temperature and pressure gradients~\cite{alkattan2018,ivanov2024}. The
extreme-conditions label applies to the initial plasma and the final collapse,
whereas the particles spend most of their time in the comparatively cold
bubble~\cite{zhang2017}. Models of this process exist and are being developed,
with bounded scope. A molecular-dynamics two-temperature model reproduces the
qualitative form of the measured size distribution and the effect of
pulse-overlap-induced porosity, but could not follow nanoparticle parameters on
the microsecond time scale and did not separate the two energy-dissipation
regimes into a rule of thumb~\cite{ivanov2024}. No model yet correlates
material and liquid properties with the primary particle diameter, and the
theoretical link between the first 100~ns and the cavitation stage is
missing~\cite{zhang2017}.

Gold is a well-studied case. Particle size and pulse fluence or pulse energy
are basic quantities, so a straightforward trend seems plausible: more energy
per pulse should give larger particles, or smaller ones once fragmentation
dominates. In the fixed corpus examined here, the reported directions do not
reduce to one trend. Articles report increases, decreases, non-monotonic
behaviour and no detectable effect, and no single direction describes the
eligible numerical series, as we report in Results. We therefore contrast the reported trends
first and then ask whether a single claim agrees with the data that its own
article supplies.

Here, we compare what articles say about the dependence of gold particle size
on pulse fluence or pulse energy with what their data show, counting directions
at the article level, including articles that state none. Agreement means
concordance under a stated trend rule, not physical confirmation.

Because the corpus is a selection rather than a representative sample, and
because pulse durations, wavelengths, fluence ranges and size measurements
differ, this heterogeneity is evidence about this corpus, not evidence that no
universal physical relationship exists or that current mechanisms fail.

\section{Related Work}
\generated[condensed from the first draft; selected evidence for the Introduction argument]{}

\subsection{Reported directions for gold.}

The primary studies cited here disagree on the direction. For femtosecond pulses at
800~nm the size increases with fluence over 60--1000~J/cm$^2$: a thermal-free
regime below about 400~J/cm$^2$ gives 3--10~nm quasi-monodisperse colloids,
whereas plasma-induced heating at higher fluence gives larger particles and a
broad distribution~\cite{kabashin2003}. For nanosecond pulses at 532~nm the
average diameter instead decreases to 5--10~nm as fluence rises to
160~J/cm$^2$~\cite{tarasenko2006}. Within one nanosecond experiment at
1064~nm, the average size decreases up to a critical fluence near
35~J/cm$^2$ and then increases, and the critical value shifts with the
medium~\cite{elsayed2013}; in another, the size is practically independent of
fluence at the fundamental wavelength and increases strongly at the second
harmonic~\cite{nikolov2011}. With pulse energy as the variable, the size grows
up to a fragmentation threshold, beyond which fragmentation
dominates~\cite{baimler2026}.

The review that records both directions in a single paragraph attributes the
contradiction to in situ fragmentation by subsequent pulses, which cannot be
excluded in batch LAL, and concludes that single-pulse fluence studies are
needed; for LAL it claims that fluence raises productivity rather than primary
particle size~\cite{zhang2017}.

\subsection{Comparability of reported sizes.}

Five analytical methods applied to the same gold colloids give
sphere-equivalent, hydrodynamic, Feret or gyration radii with different
weightings, so that reported shifts of the primary-particle mode around 10~nm
require a significance test tied to the method~\cite{letzel2018}.

\subsection{Difference from the sources.}

The five primary studies cited above are single-setup experiments, and the two
reviews report no claim-level counts and no test of a claim against the same
article's series. Neither kind of source can say how often a direction is
stated in a corpus or whether a stated direction survives comparison with the
data of the article that states it.

\section{Methods}
\PaperTODO{Describe claim extraction, series construction, and matching rules.}

\section{Verification}
\PaperTODO{Explain what the checks cover and what they do not establish.}

\section{Limitations}
\PaperTODO{Describe the conditions that limit the conclusions.}

\section{Results}
\PaperTODO{Present the main table with explicit denominators and illustrative cases.}

\section{Reproducibility}
\PaperTODO{Identify the data and code versions, environment, and run commands without revealing author identities.}

\section*{Conclusion}
\PaperTODO{State the descriptive conclusion within the scope of the corpus.}

% Mark passages with \written{...}, \generated[model]{...}, or \sourced{...}.
\begin{aiusage}
  Introduction & \generated{} &
    open-weight LLM via Yandex AI Studio &
    Cut from 569 to 413 words with the documented count script; the order of
    the narrative points compared with research/problem.md; each of the three
    citation keys (zhang2017, alkattan2018, ivanov2024) matched to its
    registry row in lit/lit.md, with DOIs not re-queried in this pass;
    quoted numbers carried over unchanged from the earlier draft and not
    re-read in this pass. \\
  Related Work & \generated{} &
    open-weight LLM via Yandex AI Studio &
    Cut from 920 to 285 words; every sentence checked against its lit/lit.md
    page record; three quoted values re-searched in the PDFs
    (kabashin2003: 400~J/cm$^2$, 3--10~nm; elsayed2013: 7--35~J/cm$^2$;
    tarasenko2006: 5--10~nm at 160~J/cm$^2$); no claim-level count written
    into prose. \\
\end{aiusage}

\begin{thebibliography}{9}

\bibitem{zhang2017}
D.~Zhang, B.~G\"okce, S.~Barcikowski,
Laser synthesis and processing of colloids: fundamentals and applications,
\emph{Chemical Reviews} 117 (2017) 3990--4103.
doi:10.1021/acs.chemrev.6b00468.

\bibitem{alkattan2018}
A.~Al-Kattan, V.~P.~Nirwan, A.~Popov, Y.~V.~Ryabchikov, G.~Tselikov, M.~Sentis,
et al.,
Recent advances in laser-ablative synthesis of bare Au and Si nanoparticles and
assessment of their prospects for tissue engineering applications,
\emph{International Journal of Molecular Sciences} 19 (2018) 1563.
doi:10.3390/ijms19061563.

\bibitem{kabashin2003}
A.~V.~Kabashin, M.~Meunier,
Synthesis of colloidal nanoparticles during femtosecond laser ablation of gold
in water,
\emph{Journal of Applied Physics} 94 (2003) 7941--7943.
doi:10.1063/1.1626793.

\bibitem{tarasenko2006}
N.~V.~Tarasenko, A.~V.~Butsen, E.~A.~Nevar, N.~A.~Savastenko,
Synthesis of nanosized particles during laser ablation of gold in water,
\emph{Applied Surface Science} 252 (2006) 4439--4444.
doi:10.1016/j.apsusc.2005.07.150.

\bibitem{elsayed2013}
K.~A.~Elsayed, H.~Imam, M.~A.~Ahmed, R.~Ramadan,
Effect of focusing conditions and laser parameters on the fabrication of gold
nanoparticles via laser ablation in liquid,
\emph{Optics \& Laser Technology} 45 (2013) 495--502.
doi:10.1016/j.optlastec.2012.06.004.

\bibitem{nikolov2011}
A.~S.~Nikolov, N.~N.~Nedyalkov, R.~G.~Nikov, P.~A.~Atanasov, M.~T.~Alexandrov,
Characterization of Ag and Au nanoparticles created by nanosecond pulsed laser
ablation in double distilled water,
\emph{Applied Surface Science} 257 (2011) 5278--5282.
doi:10.1016/j.apsusc.2010.10.146.

\bibitem{baimler2026}
I.~V.~Baimler, I.~A.~Popov, A.~V.~Simakin, S.~V.~Gudkov,
A simple three-step method for the synthesis of submicron gold particles: the
influence of laser irradiation duration, pulse energy, laser pulse duration,
and initial concentration of nanoparticles in the colloid,
\emph{Nanomaterials} 16 (2026) 79.
doi:10.3390/nano16020079.

\bibitem{ivanov2024}
D.~S.~Ivanov, P.~Shakhov, G.~Tikhonowsky, A.~A.~Popov, A.~N.~Mayorov,
I.~N.~Zavestovskaya, et al.,
Ultrafast laser ablation of gold in liquids: effect of laser pulse
overlap-induced surface porosity on size distribution of formed nanoparticles,
\emph{Applied Surface Science} 643 (2024) 158662.
doi:10.1016/j.apsusc.2023.158662.

\bibitem{letzel2018}
A.~Letzel, B.~G\"okce, A.~Menzel, A.~Plech, S.~Barcikowski,
Primary particle diameter differentiation and bimodality identification by
five analytical methods using gold nanoparticle size distributions synthesized
by pulsed laser ablation in liquids,
\emph{Applied Surface Science} 435 (2018) 743--751.
doi:10.1016/j.apsusc.2017.11.130.

\end{thebibliography}

\end{document}
